\documentclass[11pt]{article}
\usepackage[T1]{fontenc}
\usepackage[utf8]{inputenc}
\usepackage{lmodern}
\usepackage[margin=1in]{geometry}
\usepackage{amsmath,amssymb}
\usepackage{booktabs}
\usepackage[numbers,sort]{natbib}
\usepackage{url}
\usepackage[hidelinks]{hyperref}

\title{Rejecting Explosion Is Not Paraconsistency:\\How Small Language Models Reason from Inconsistent Premises}
% TODO(author): replace with your name and affiliation before submission.
\author{Author Name\\Affiliation\\\texttt{email@example.org}}
\date{Preprint, November 2026}

\begin{document}
\maketitle

\begin{abstract}
Logic benchmarks for language models score answers against classical logic, so a model that accepts
disjunctive syllogism (DS) but rejects explosion ($A, \neg A \vdash B$) counts as making a classical error,
or as hinting at a paraconsistent logic. In a pilot on seven local open-weight models (7--30B parameters)
every model showed exactly this pattern, which fits none of six candidate logics (classical,
intuitionistic, K3, LP, FDE, orthomodular quantum logic). We ran four follow-up experiments, each with
predictions committed to version control before scoring. The rejection of explosion survives conclusions
that are on topic and an explicit definition of validity. It is not a consequence relation of any of the
six logics, and not a refusal to reason from inconsistent premises at all: the models still conclude
$C$ from $A, \neg A, C$. What they refuse is any conclusion that rests on the contradicted sentence,
including $A$ itself from $A, \neg A$, which every candidate logic validates. For the one model that
passed every reliability check (qwen2.5:7b-instruct), a pre-specified hypothesis based on maximal
consistent subsets (Rescher and Manor) was refuted, and the model's profile instead matched a
pre-specified rival exactly: drop what the premises say about the contradicted atom and reason
classically with the rest. A formalisation of that rival, fixed after the data, agrees with all 14
schemata we measured for this model. We argue that explosion items in logic benchmarks measure a
strategy for handling inconsistency, not adherence to a logic, and we report the models and checks that
did not support these conclusions.
\end{abstract}

\section{Introduction}

Putnam held that the choice of logic is, in principle, open to revision on empirical grounds
\citep{putnam1968}, and later that at least one principle, ``not every statement is both true and
false'', survives every rational revision \citep{putnam1978}. Language models invite a version of the
first question that can actually be measured: a model's inferential behaviour is learned from text, so
one can ask which logic, if any, describes it.

We built a battery of inference schemata whose verdicts differ across six candidate logics and asked
local models whether each argument is valid. The first result was a pattern that none of the six logics
produces. All seven pilot models accepted DS ($A \lor B, \neg A \vdash B$) and rejected explosion
($A, \neg A \vdash B$). Classical, intuitionistic and K3 logic validate both; LP and FDE validate
neither; orthomodular quantum logic validates explosion and rejects DS. The pattern is also puzzling for
a second reason. C.~I.~Lewis's derivation of explosion \citep{lewis1932} uses only $\lor$-introduction
($A \vdash A \lor B$), DS, and chaining. A reasoner that accepts both steps but rejects the conclusion
either has a non-transitive consequence relation or is not applying a consequence relation at all.

This paper reports four follow-up experiments that decide between these readings. Our contributions:
\begin{itemize}
\item The rejection of explosion is robust across five model families and does not go away when the
conclusion is on topic or when the prompt states the classical definition of validity
(Section~\ref{sec:pragmatics}).
\item It is not evidence for paraconsistency or for non-transitive consequence: models also reject
$A, \neg A \vdash A$, valid in every candidate logic, while still accepting $A, \neg A, C \vdash C$
(Sections~\ref{sec:chain}--\ref{sec:inconsistent}).
\item For the confirmatory model, a pre-specified hypothesis based on maximal consistent subsets was
refuted, and a pre-specified rival, ``discard the contradicted atom'', matched exactly
(Section~\ref{sec:mcs}).
\item All items, scores, analysis code and the dated record of predictions are released.
\end{itemize}

\section{Related work}

Benchmarks for logical reasoning in language models score answers against classical logic: rule
chaining over synthetic theories \citep{clark2020,saparov2023}, first-order reasoning with
prover-checked annotations \citep{han2024}, single inference rules across propositional, first-order and
non-monotonic reasoning \citep{parmar2024}, and multi-step combinations of rules \citep{patel2024}. They
report accuracy against one target. Our question differs: we ask which logic, if any, fits the
answers, and what a systematic departure from classical logic is.
\citet{lampinen2024} show that model validity judgements, like human ones, depend on whether the
conclusion is believable. We control content in a similar spirit, but our manipulation targets relevance
and inconsistency rather than belief.

The candidate logics are classical logic (CL), intuitionistic logic (IL), strong Kleene logic K3
\citep{kleene1952}, Priest's LP \citep{priest1979}, first-degree entailment FDE \citep{dunn1976}, and
orthomodular quantum logic QL \citep{birkhoff1936}, the logic behind Putnam's 1968 proposal. For
inconsistent premise sets we also consider consequence from maximal consistent subsets
\citep{rescher1970}.

\section{Method}

\paragraph{Items.}
Each schema (premises and conclusion over atoms $A, B, C$) is instantiated with everyday sentences
(``the shop is open''), nonsense sentences (``the blicket is wuggy'') and symbolic notation ($P \land
\neg Q$), and asked in three framings: ``Does the following conclusion follow from the premises?'',
``Suppose the following are all true\ldots{} Must it then be true that\ldots?'', and ``Is this argument
logically valid?''. Every prompt ends with ``Answer yes or no.''. Gold verdicts are computed, not
hand-labelled: exhaustive truth tables for CL, K3, LP and FDE, a search over Kripke models with up to
three worlds for IL, and valuations into the smallest non-distributive orthomodular lattice for QL. A
test file checks the computed verdicts against hand-derived ones.

\paragraph{Scoring.}
For each item we take the first-token log-probabilities from the model (temperature 0, top 20
tokens) and compute $P(\text{yes}) = m_{\text{yes}} / (m_{\text{yes}} + m_{\text{no}})$, where $m$ is the
probability mass on tokens that read ``yes'' or ``no''. Each item is scored twice, with ``Answer yes or
no'' and ``Answer no or yes'', and the two values are averaged. A cell is read as acceptance when its mean
$P(\text{yes}) > 0.5$.

\paragraph{Models.}
All models were run locally through Ollama. The pilot used qwen2.5:7b-instruct, qwen2.5:14b,
qwen3.5:9b, qwen3.8:27b (both with reasoning disabled), qwen3-coder:30b, hermes3:8b and llama3.1:latest. The follow-ups used
qwen2.5:7b-instruct, gemma3:12b, hermes3:8b and mistral:7b; llama3.1 was dropped after the first
follow-up (below). Model digests are listed in Appendix~\ref{app:digests}.

\paragraph{Reliability rule.}
Fixed before the follow-ups: a model is ``unstable'' in a run if its mean $P(\text{yes})$ on the two
fallacies (affirming the consequent, denying the antecedent) exceeds 0.25, or if more than 25\% of items
put less than half of the first-token mass on yes/no in either answer order. Conclusions about an
unstable model are reported as exploratory. Runs without fallacy items are judged on answer mass only.

\paragraph{Record of predictions.}
Thresholds and the primary analyses were frozen in a preregistration file before the first follow-up was
scored; each follow-up's schemata, item file and predicted response patterns were committed before any
of its scores (commit history in the released repository; Appendix~\ref{app:timeline}). The pilot was
analysed post hoc and is exploratory. One item set was changed before scoring, and we report it: the
first version of the ``formal'' framing put the definition of validity first. In a smoke test
qwen2.5:7b restated the definition as a question instead of answering (yes/no mass 0.03 and 0.09 on the
formal item measured, and the same restatement on eight further formal items). We compared three
rewordings on 12 items by answer mass alone and used the one that put all mass on yes/no in all 24
probes.

% Generated by paper/make_tables.py; do not edit by hand.
\begin{table}[t]\centering\small
\caption{Pilot (exploratory): mean $P(\text{yes})$ on disjunctive syllogism (DS) and explosion over all content types and framings, yes-rate on the two fallacies, share of items with yes/no first-token mass $<0.5$, and the reliability verdict.}
\label{tab:pilot}
\begin{tabular}{lrrrrl}\toprule
Model & DS & Explosion & Fallacy yes & Low mass & Reliability \\\midrule
hermes3:8b & 0.86 & 0.01 & 0.34 & 0\% & unstable \\
llama3.1:latest & 0.61 & 0.28 & 0.49 & 62\% & unstable \\
qwen2.5:14b & 1.00 & 0.00 & 0.21 & 0\% & ok \\
qwen2.5:7b-instruct & 1.00 & 0.00 & 0.08 & 18\% & ok \\
qwen3-coder:30b & 1.00 & 0.00 & 0.16 & 6\% & ok \\
qwen3.5:9b & 0.99 & 0.13 & 0.11 & 0\% & ok \\
qwen3.8:27b & 1.00 & 0.07 & 0.02 & 22\% & ok \\
\bottomrule\end{tabular}\end{table}

\begin{table}[t]\centering\small
\caption{Pragmatic controls: mean $P(\text{yes})$ on explosion with unrelated (everyday) and topically related atoms under the three default framings, and under the framing that states the classical definition of validity (both contents). Fallacy yes-rate, low-mass share and reliability are over the whole run.}
\label{tab:pragmatics}
\begin{tabular}{lrrrrrl}\toprule
Model & Unrelated & Related & Formal & Fallacy yes & Low mass & Reliability \\\midrule
qwen2.5:7b-instruct & 0.00 & 0.00 & 0.00 & 0.14 & 8\% & ok \\
gemma3:12b & 0.02 & 0.00 & 0.02 & 0.41 & 6\% & unstable \\
hermes3:8b & 0.00 & 0.00 & 0.03 & 0.28 & 0\% & unstable \\
mistral:7b & 0.08 & 0.11 & 0.11 & 0.45 & 8\% & unstable \\
llama3.1:latest & 0.18 & 0.23 & 0.39 & 0.46 & 52\% & unstable \\
\bottomrule\end{tabular}\end{table}

\begin{table*}[t]\centering\small
\caption{Profile over the follow-up runs (pooled across runs where a schema recurs). Left: verdicts of classical logic, LP, universal maximal-consistent-subset consequence (MCS-U) and MCS-U after splitting conjunctive premises (Split, formalised post hoc). Right: mean $P(\text{yes})$; for the confirmatory model qwen2.5:7b, with a 95\% instance-bootstrap interval. Upper block: consistent premises; lower block: inconsistent premises. Every form is classically valid; forms marked $\times$ under LP are invalid in LP and FDE, and all others are valid in all six candidate logics (IL up to three worlds, QL in MO2). Bottom: number of schemata whose binarised mean disagrees with each candidate.}
\label{tab:profile}
\begin{tabular}{lccccrrrr}\toprule
Form & CL & LP & MCS-U & Split & qwen2.5:7b & gemma3 & hermes3 & mistral \\\midrule
$A, A\to B \vdash B$ & \checkmark & $\times$ & \checkmark & \checkmark & 0.95 [0.85, 1.00] & 0.97 & 0.55 & 0.65 \\
$A\lor B, \neg A \vdash B$ & \checkmark & $\times$ & \checkmark & \checkmark & 1.00 [1.00, 1.00] & 1.00 & 0.83 & 0.67 \\
$A\lor B, \neg A, \neg B\lor C \vdash C$ & \checkmark & $\times$ & \checkmark & \checkmark & 0.85 [0.77, 0.92] & 1.00 & 0.54 & 0.38 \\
$A \vdash A\lor B$ & \checkmark & \checkmark & \checkmark & \checkmark & 0.84 [0.77, 0.91] & 0.94 & 0.41 & 0.40 \\
$A\land C \vdash C$ & \checkmark & \checkmark & \checkmark & \checkmark & 0.99 [0.96, 1.00] & 0.99 & 0.80 & 0.67 \\
$A, C \vdash C$ & \checkmark & \checkmark & \checkmark & \checkmark & 0.83 [0.74, 0.92] & 1.00 & 0.70 & 0.69 \\
\midrule
$A, \neg A \vdash B$ & \checkmark & $\times$ & $\times$ & $\times$ & 0.00 [0.00, 0.00] & 0.08 & 0.02 & 0.11 \\
$A, \neg A \vdash A$ & \checkmark & \checkmark & $\times$ & $\times$ & 0.01 [0.00, 0.02] & 0.24 & 0.08 & 0.03 \\
$A, \neg A \vdash A\land\neg A$ & \checkmark & \checkmark & $\times$ & $\times$ & 0.00 [0.00, 0.01] & 0.39 & 0.11 & 0.12 \\
$A, \neg A \vdash A\lor B$ & \checkmark & \checkmark & $\times$ & $\times$ & 0.19 [0.10, 0.29] & 0.22 & 0.16 & 0.13 \\
$A, \neg A, A\lor B \vdash B$ & \checkmark & $\times$ & $\times$ & $\times$ & 0.32 [0.27, 0.38] & 0.91 & 0.27 & 0.09 \\
$A, \neg A, C \vdash C$ & \checkmark & \checkmark & \checkmark & \checkmark & 0.76 [0.69, 0.82] & 1.00 & 0.49 & 0.46 \\
$A\land C, C\land\neg A \vdash C$ & \checkmark & \checkmark & \checkmark & \checkmark & 0.87 [0.79, 0.94] & 0.98 & 0.69 & 0.62 \\
$A\land C, \neg A \vdash C$ & \checkmark & \checkmark & $\times$ & \checkmark & 0.94 [0.89, 0.98] & 0.80 & 0.39 & 0.42 \\
\midrule
\multicolumn{5}{r}{mismatches with CL} & 5 & 4 & 8 & 9 \\
\multicolumn{5}{r}{mismatches with LP} & 6 & 7 & 9 & 8 \\
\multicolumn{5}{r}{mismatches with MCS-U} & 1 & 2 & 2 & 3 \\
\multicolumn{5}{r}{mismatches with Split} & 0 & 1 & 3 & 4 \\
\bottomrule\end{tabular}\end{table*}


\section{Results}

\subsection{Pilot (exploratory)}
Table~\ref{tab:pilot} shows the pattern that started the study. DS is accepted (mean
$P(\text{yes})$ 0.61--1.00) and explosion rejected (0.00--0.28) by all seven models. Five of seven pass
the reliability rule. On the six conditional-free schemata the fit to each logic is uninformative about
this pattern: CL is penalised for explosion and LP for DS, so the two tie.

\subsection{Is the rejection pragmatic?}\label{sec:pragmatics}
Two cheap explanations would dissolve the puzzle. The model may be judging relevance (``the oven is hot
and not hot, so the lights are off'' is a non sequitur), or it may read ``follows'' conversationally. We
re-asked explosion with atoms that share a subject (``the kettle is full and not full, so the kettle is
hot'') and with a framing that states that an argument is valid if and only if no possible situation makes
all premises true and the conclusion false. Neither manipulation moved any model toward acceptance
(Table~\ref{tab:pragmatics}); the largest mean is 0.11 for the models we kept. The formal framing did not
raise the yes-rate on fallacies for qwen2.5:7b (0.09 and 0.31). Four models failed the fallacy check in
this run, almost entirely through denying the antecedent, which suggests a biconditional reading of
``if''. We did not change the threshold after seeing this; those models are exploratory from here on.
llama3.1 also put less than half of the mass on yes/no in 52\% of items and was dropped.

\subsection{Which step of Lewis's argument is refused?}\label{sec:chain}
We decomposed the derivation: $\lor$-introduction, DS with the inconsistent pair present
($A, \neg A, A\lor B \vdash B$), two chained DS steps without a contradiction, and, as a check on whether
anything at all is drawn from inconsistent premises, $A, \neg A \vdash A$. Four hypotheses with written
predictions were written on 14 October and committed before any of these items were scored.
Table~\ref{tab:profile} shows the result in the rows for DS, chained DS, $\lor$-introduction, explosion,
$A, \neg A \vdash A$ and the DS lemma. qwen2.5:7b accepts $\lor$-introduction (0.84), DS (1.00) and
chained DS (0.85) and rejects the DS lemma with a contradiction present (0.32). It also rejects
$A, \neg A \vdash A$ (0.01), which every candidate logic validates. This rules out a non-transitive
consequence relation for this model: it does not accept every step and refuse the conclusion; it refuses
the step that runs through the contradiction. gemma3:12b is different: it accepts the DS lemma (0.91)
but not explosion, which is the non-transitive profile, apart from its rejection of $A, \neg A \vdash A$.

\subsection{Global aversion or a local filter?}\label{sec:inconsistent}
A ``no'' to $A, \neg A \vdash A$ has two readings. Either the model draws nothing from an inconsistent
premise set, or it treats the premises as conflicting evidence about $A$ and declines to conclude
anything about $A$. A third schema separates them: $A, \neg A, C \vdash C$. qwen2.5:7b accepts it (0.76;
0.88--0.98 in natural language, 0.39--0.47 in symbolic notation) and so does gemma3:12b (1.00), while both
reject restating the contradiction itself ($A, \neg A \vdash A \land \neg A$: 0.00 and 0.39). By the
reading committed before this run, the ``global aversion'' diagnosis is withdrawn for both models.
hermes3:8b and mistral:7b sit just under 0.5 on the bystander item (0.49, 0.46) and are weak on
restating a premise even without a contradiction.

\subsection{A pre-specified test: maximal consistent subsets}\label{sec:mcs}
After the previous run, qwen2.5:7b's profile over eight schemata matched consequence from maximal
consistent subsets (MCS): a conclusion is accepted if it follows classically from every maximal
consistent subset of the premises \citep{rescher1970}. Because that hypothesis was formed after seeing
the data, we tested it on new items chosen to separate it from four rivals, with the confirmatory model,
decision rule and an inconclusive band (0.40--0.60 on the deciding item) fixed in advance. The run
had seven schemata: four new ones ($A \land C \vdash C$, $A \land C, C \land \neg A \vdash C$,
$A \land C, \neg A \vdash C$ and $A, \neg A \vdash A \lor B$) and three anchors with new instances. The rivals were: consequence from the premises outside every minimal
inconsistent subset; consequence from at least one MCS; global aversion; and ``discard the contradicted
atom'', stated informally as: drop what the premises say about an atom they both assert and deny, keep
everything else, including the other conjunct of a conjunction.

The deciding item for the last rival is $A \land C, \neg A \vdash C$: MCS rejects it, since $\{\neg A\}$
is a maximal consistent subset that does not entail $C$. qwen2.5:7b accepts it (0.94, 95\% interval
0.89--0.98). The MCS hypothesis is therefore refuted for the confirmatory model, and the model's
seven-item profile matches ``discard the contradicted atom'' exactly and uniquely. Among the exploratory
models, hermes3:8b matched MCS exactly in this run (deciding item 0.39), gemma3:12b matched no candidate
(it again accepts the DS lemma), and mistral:7b matched none.

\paragraph{A formalisation, fixed after the data.}
``Discard the contradicted atom'' coincides, on our items, with MCS consequence computed after splitting
each conjunctive premise into its conjuncts. Across all 14 schemata in Table~\ref{tab:profile} it
disagrees with qwen2.5:7b on none, against one for plain MCS and five for CL. The relation is
syntax-sensitive (it treats $A \land C$ differently from the premise pair $A, C$ only before splitting),
so it is not a consequence relation on propositions. We report it as a description to be tested, not as
a finding.

\section{Discussion}

\paragraph{What explosion items measure.}
A benchmark that counts the rejection of explosion as a classical error, or as a sign of a paraconsistent
logic, misreads it. The models we studied do not reject explosion because they lack $\lor$-introduction
or DS, and not because they hold a logic in which contradictions are harmless. They reject it because
they will not draw conclusions that depend on a sentence the premises both assert and deny, and an
explosion item has no other support for its conclusion. The same strategy makes them reject inferences
that every candidate logic accepts, such as $A, \neg A \vdash A$.

\paragraph{Relation to Putnam.}
qwen2.5:7b never concludes $A \land \neg A$ from $A$ and $\neg A$ (mean 0.00). One can read this as a
behavioural counterpart of Putnam's minimal principle of contradiction \citep{putnam1978}: the model
does not treat any statement as both true and false. The analogy is limited. Putnam's principle is
normative and about rational revision; ours is a regularity in a 7B model's answers, and a regularity
says nothing about whether giving it up could be rational.

\paragraph{Heterogeneity.}
The models do not share one strategy. gemma3:12b runs DS through a contradiction (0.91) but still
rejects explosion, a profile none of our candidates predicts. The two weaker models are close to
refusing everything from inconsistent premises. We do not have a single account that covers all four.

\section{Limitations}

\begin{itemize}
\item Only qwen2.5:7b-instruct passed every reliability check; the other follow-up models failed the
fallacy check in the run that had fallacy items, and their profiles are exploratory. All follow-up models
are local, open-weight and 7--12B parameters. No closed API model or larger model was tested in the
follow-ups.
\item The final formalisation (MCS after splitting conjunctions) was written after seeing the data and
has not been tested on items where it and its rivals diverge, such as premises that split only after De
Morgan's laws.
\item First-token probabilities measure the model's immediate answer, not what it would conclude after
reasoning. Items are English, template-generated and propositional.
\item The reliability rule's fallacy check is driven mostly by denying the antecedent, a conditional
schema whose natural-language reading is contested. We kept the rule as frozen.
\item The record of predictions is a git history; commit timestamps are set by the author's machine.
% TODO(author): make the repository public or archive it (e.g. Zenodo) before submission and cite the DOI here.
\end{itemize}

\section*{Code and data availability}
Items, score files, model digests, analysis code (\texttt{battery/}), the per-run notes with predictions,
and the script that computes every number in this paper (\texttt{paper/make\_tables.py}) are in the
project repository.
% TODO(author): add the public repository URL or archive DOI.

% TODO(author): many venues expect a statement on AI assistance; edit or remove.
\section*{Use of AI tools}
Code, analyses and a first draft of this text were prepared with the help of an AI coding assistant.
The author checked the analyses and is responsible for the content.

\bibliographystyle{plainnat}
\bibliography{refs}

\appendix
\section{Timeline of committed predictions}\label{app:timeline}
\begin{center}\small
\begin{tabular}{lll}\toprule
Run & Items and predictions committed & Scores committed \\\midrule
Explosion chain & 2026-10-14 (notes); repository snapshot 2026-10-25 & 2026-11-24 18:57 UTC \\
Thresholds frozen & 2026-11-24 16:11 UTC & -- \\
Pragmatic controls & 2026-11-24 17:00 UTC (reworded 17:18) & 2026-11-24 18:57 UTC \\
Inconsistent premises & 2026-11-24 19:04 UTC & 2026-11-24 19:57 UTC \\
MCS test & 2026-11-24 20:03 UTC & 2026-11-24 20:12 UTC \\
\bottomrule\end{tabular}
\end{center}
The explosion-chain hypotheses were written on 2026-10-14; the analysis thresholds applied to that run
were frozen on 2026-11-24, before it was scored. qwen2.5:7b had been scored on the pragmatic-control run
before the other models.

\section{Model digests}\label{app:digests}
% Generated by paper/make_tables.py; do not edit by hand.
\begin{table}[h]\centering\small
\caption{Ollama model digests (first 12 hex digits) recorded for the follow-up runs.}
\label{tab:digests}
\begin{tabular}{ll}\toprule Model & Digest \\\midrule
qwen2.5:7b-instruct & \texttt{845dbda0ea48} \\
gemma3:12b & \texttt{b92b218582c5} \\
hermes3:8b & \texttt{4f6b83f30b62} \\
mistral:7b & \texttt{6577803aa9a0} \\
\bottomrule\end{tabular}\end{table}


\end{document}
